
\documentclass{article}
\usepackage{colortbl}
\usepackage{makecell}
\usepackage{multirow}
\usepackage{supertabular}

\begin{document}

\newcounter{utterance}

\centering \large Interaction Transcript for game `clean\_up', experiment `1\_medium\_5obj\_sl', episode 2 with gemma{-}4{-}26b{-}a4b{-}it.
\vspace{24pt}

{ \footnotesize  \setcounter{utterance}{1}
\setlength{\tabcolsep}{0pt}
\begin{supertabular}{c@{$\;$}|p{.15\linewidth}@{}p{.15\linewidth}p{.15\linewidth}p{.15\linewidth}p{.15\linewidth}p{.15\linewidth}}
    \# & $\;$A & \multicolumn{4}{c}{Game Master} & $\;\:$B\\
    \hline

    \theutterance \stepcounter{utterance}  
    & & \multicolumn{4}{p{0.6\linewidth}}{
        \cellcolor[rgb]{0.9,0.9,0.9}{
            \makecell[{{p{\linewidth}}}]{
                \texttt{\tiny{[P1$\langle$GM]}}
                \texttt{Jaz sem vaš vodja igre, vi pa igrate skupinsko igro z naslednjo mrežo kot igralno ploščo:} \\
\\ 
\texttt{```} \\
\texttt{ 1234567} \\
\texttt{╔═╤═╤═══╗} \\
\texttt{║◌│◌│◌◌◌║ 1} \\
\texttt{╟─┴─┼─┬─╢ 2} \\
\texttt{║◌◌◌│◌│T║ 3} \\
\texttt{║◌H◌├─┘◌║ 4} \\
\texttt{║◌┌─┘◌◌W║ 5} \\
\texttt{╟─┤◌◌◌◌◌║ 6} \\
\texttt{║I│◌C◌◌◌║ 7} \\
\texttt{╚═╧═════╝} \\
\\ 
\texttt{```} \\
\\ 
\texttt{* Na zgornjem robu so prikazane koordinate x, ki naraščajo proti desni, na desnem robu pa koordinate y, ki naraščajo navzdol.} \\
\texttt{* Na mrežo so naključno postavljeni naslednji predmeti: 'W', 'I', 'T', 'C', 'H'.} \\
\\ 
\\ 
\texttt{Drugi igralec vidi različico igralne plošče, na kateri so predmeti postavljeni na različna naključna mesta. Vi ne morete videti plošče drugega igralca, on pa ne more videti vaše.} \\
\\ 
\\ 
\texttt{**Cilj:**} \\
\\ 
\texttt{Oba igralca morata premikati predmete na svojem ozadju tako, da se enaki predmeti znajdejo na istih koordinatah. Z drugim igralcem se morate dogovoriti o skupni ciljni konfiguraciji.} \\
\\ 
\\ 
\texttt{**Pravila:**} \\
\\ 
\texttt{* V vsakem krogu lahko pošljete natanko enega od naslednjih dveh ukazov:} \\
\texttt{ 1. `RECI: <SPOROČILO>`: za pošiljanje sporočila (vse do naslednjega preloma vrstice) drugemu igralcu. Posredoval ga bom vašemu partnerju.} \\
\texttt{2. `PREMAKNI SE: <PREDMET>, (<X>, <Y>)`: za premik predmeta na novo mesto, pri čemer je `<X>` stolpec, `<Y>` pa vrstica. Obvestil vas bom, ali je bila vaša premestitev uspešna ali ne.} \\
\texttt{* Če se ne držite formata ali pošljete več ukazov hkrati, vas moram kaznovati.} \\
\texttt{* Če oba igralca zbereta več kot 12 kazni, oba izgubita igro.} \\
\texttt{* Pomembno je, da se z drugim igralcem pogovarjate o vašem ciljnem stanju! Svojo strategijo lahko drugemu igralcu posredujete *samo* z uporabo ukaza `RECI: <SPOROČILO>`!} \\
\\ 
\\ 
\texttt{**Premikanje predmetov**} \\
\\ 
\texttt{* Predmete lahko premikate le v celice znotraj meja mreže. Ciljna celica mora biti prazna, tj. vsebovati mora samo simbol '◌'.} \\
\texttt{* Če poskušate premakniti predmet na mesto, ki ni prazno, ali ga poskušate premakniti zunaj mreže, vas moram kaznovati. Dobili boste še en poskus.} \\
\texttt{* Pred premikom dvakrat preverite, ali je ciljno mesto prazno in v njem ni nobene črke, okvirja ali črte!} \\
\\ 
\\ 
\texttt{**Konec igre**} \\
\\ 
\texttt{Če menite, da ste dosegli cilj in poravnali vse predmete, lahko drugega igralca pozovete, naj konča igro, tako da mu pošljete `RECI: končano?`. Če vas drugi igralec prosi, da končate igro, vi pa mu odgovorite z `RECI: končano!`, se igra konča.} \\
\\ 
\texttt{Oba igralca zmagata, če se igra konča v 20 krogih, pri čemer je en krog definiran tako, da dva igralca pošljeta vsak po en veljaven ukaz.} \\
\\ 
\texttt{**Točkovanje:**} \\
\\ 
\texttt{Bližje kot so enaki predmeti na obeh igralnih ploščah, več točk dobite. Kazni zmanjšajo število točk. Ali lahko premagate rekord?} \\
\\ 
\\ 
\texttt{Pošljite sporočilo drugemu igralcu, da začnete igro!} \\
            }
        }
    }
    & \\ \\

    \theutterance \stepcounter{utterance}  
    & \multicolumn{4}{p{0.6\linewidth}}{
        \cellcolor[rgb]{0.9,0.9,0.9}{
            \makecell[{{p{\linewidth}}}]{
                \texttt{\tiny{[P1$\rangle$GM]}}
                \texttt{RECI: Pozdravljeni! Začnimo igro. Najprej se morava dogovoriti o ciljni konfiguraciji. Ker ne vem, kje so predmeti pri vas, predlagam, da poskusiva postaviti vse predmete v določene fiksne koordinate.} \\
\\ 
\texttt{나predlagam naslednjo ciljno konfiguracijo:} \\
\texttt{W: (1, 1)} \\
\texttt{I: (2, 1)} \\
\texttt{T: (3, 1)} \\
\texttt{C: (4, 1)} \\
\texttt{H: (5, 1)} \\
\\ 
\texttt{Ali se strinjaš na to? Če ne, predlaga drugo sestavo.} \\
            }
        }
    }
    & & \\ \\

    \theutterance \stepcounter{utterance}  
    & & & \multicolumn{2}{p{0.3\linewidth}}{
        \cellcolor[rgb]{0.9,0.9,0.9}{
            \makecell[{{p{\linewidth}}}]{
                \texttt{\tiny{[GM$|$GM]}}
                \texttt{RECI: Pozdravljeni! Začnimo igro. Najprej se morava dogovoriti o ciljni konfiguraciji. Ker ne vem, kje so predmeti pri vas, predlagam, da poskusiva postaviti vse predmete v določene fiksne koordinate.} \\
\\ 
\texttt{나predlagam naslednjo ciljno konfiguracijo:} \\
\texttt{W: (1, 1)} \\
\texttt{I: (2, 1)} \\
\texttt{T: (3, 1)} \\
\texttt{C: (4, 1)} \\
\texttt{H: (5, 1)} \\
\\ 
\texttt{Ali se strinjaš na to? Če ne, predlaga drugo sestavo.} \\
            }
        }
    }
    & & \\ \\

    \theutterance \stepcounter{utterance}  
    & & & \multicolumn{2}{p{0.3\linewidth}}{
        \cellcolor[rgb]{0.9,0.9,0.9}{
            \makecell[{{p{\linewidth}}}]{
                \texttt{\tiny{[GM$|$GM]}}
                \texttt{Invalid format: tail is not empty: '} \\
\\ 
\texttt{나predlagam naslednjo ciljno konfiguracijo:} \\
\texttt{W: (1, 1)} \\
\texttt{I: (2, 1)} \\
\texttt{T: (3, 1)} \\
\texttt{C: (4, 1)} \\
\texttt{H: (5, 1)} \\
\\ 
\texttt{Ali se strinjaš na to? Če ne, predlaga drugo sestavo.} \\
            }
        }
    }
    & & \\ \\

    \theutterance \stepcounter{utterance}  
    & & \multicolumn{4}{p{0.6\linewidth}}{
        \cellcolor[rgb]{0.9,0.9,0.9}{
            \makecell[{{p{\linewidth}}}]{
                \texttt{\tiny{[P1$\langle$GM]}}
                \texttt{Kazen: Vaše sporočilo ne sme vsebovati ničesar po ukazu!} \\
\\ 
\texttt{Prepričajte se, da vaš odgovor vsebuje samo `RECI: <SPOROČILO>` ali `PREMAKNI SE: <PREDMET>, (<X>, <Y>)` in nič drugega!} \\
\\ 
\texttt{Skupaj ste zbrali 1 od 12 kazni.} \\
\texttt{Prosim, poskusite znova!} \\
            }
        }
    }
    & \\ \\

    \theutterance \stepcounter{utterance}  
    & \multicolumn{4}{p{0.6\linewidth}}{
        \cellcolor[rgb]{0.9,0.9,0.9}{
            \makecell[{{p{\linewidth}}}]{
                \texttt{\tiny{[P1$\rangle$GM]}}
                \texttt{RECI: Pozdravljeni! Dogovoriva se o naslednji ciljni konfiguraciji: W:(1,1), I:(2,1), T:(3,1), C:(4,1), H:(5,1). Ali se strinjaš?} \\
            }
        }
    }
    & & \\ \\

    \theutterance \stepcounter{utterance}  
    & & & \multicolumn{2}{p{0.3\linewidth}}{
        \cellcolor[rgb]{0.9,0.9,0.9}{
            \makecell[{{p{\linewidth}}}]{
                \texttt{\tiny{[GM$|$GM]}}
                \texttt{RECI: Pozdravljeni! Dogovoriva se o naslednji ciljni konfiguraciji: W:(1,1), I:(2,1), T:(3,1), C:(4,1), H:(5,1). Ali se strinjaš?} \\
            }
        }
    }
    & & \\ \\

    \theutterance \stepcounter{utterance}  
    & & \multicolumn{4}{p{0.6\linewidth}}{
        \cellcolor[rgb]{0.9,0.9,0.9}{
            \makecell[{{p{\linewidth}}}]{
                \texttt{\tiny{[GM$\rangle$P2]}}
                \texttt{Jaz sem vaš vodja igre, vi pa igrate skupinsko igro z naslednjo mrežo kot igralno ploščo:} \\
\\ 
\texttt{```} \\
\texttt{ 1234567} \\
\texttt{╔═╤═╤═══╗} \\
\texttt{║◌│◌│◌◌◌║ 1} \\
\texttt{╟─┴─┼─┬─╢ 2} \\
\texttt{║◌◌◌│◌│◌║ 3} \\
\texttt{║◌◌C├─┘◌║ 4} \\
\texttt{║◌┌─┘◌T◌║ 5} \\
\texttt{╟─┤◌W◌◌H║ 6} \\
\texttt{║◌│◌I◌◌◌║ 7} \\
\texttt{╚═╧═════╝} \\
\\ 
\texttt{```} \\
\\ 
\texttt{* Na zgornjem robu so prikazane koordinate x, ki naraščajo proti desni, na desnem robu pa koordinate y, ki naraščajo navzdol.} \\
\texttt{* Na mrežo so naključno postavljeni naslednji predmeti: 'W', 'I', 'T', 'C', 'H'.} \\
\\ 
\\ 
\texttt{Drugi igralec vidi različico igralne plošče, na kateri so predmeti postavljeni na različna naključna mesta. Vi ne morete videti plošče drugega igralca, on pa ne more videti vaše.} \\
\\ 
\\ 
\texttt{**Cilj:**} \\
\\ 
\texttt{Oba igralca morata premikati predmete na svojem ozadju tako, da se enaki predmeti znajdejo na istih koordinatah. Z drugim igralcem se morate dogovoriti o skupni ciljni konfiguraciji.} \\
\\ 
\\ 
\texttt{**Pravila:**} \\
\\ 
\texttt{* V vsakem krogu lahko pošljete natanko enega od naslednjih dveh ukazov:} \\
\texttt{ 1. `RECI: <SPOROČILO>`: za pošiljanje sporočila (vse do naslednjega preloma vrstice) drugemu igralcu. Posredoval ga bom vašemu partnerju.} \\
\texttt{2. `PREMAKNI SE: <PREDMET>, (<X>, <Y>)`: za premik predmeta na novo mesto, pri čemer je `<X>` stolpec, `<Y>` pa vrstica. Obvestil vas bom, ali je bila vaša premestitev uspešna ali ne.} \\
\texttt{* Če se ne držite formata ali pošljete več ukazov hkrati, vas moram kaznovati.} \\
\texttt{* Če oba igralca zbereta več kot 12 kazni, oba izgubita igro.} \\
\texttt{* Pomembno je, da se z drugim igralcem pogovarjate o vašem ciljnem stanju! Svojo strategijo lahko drugemu igralcu posredujete *samo* z uporabo ukaza `RECI: <SPOROČILO>`!} \\
\\ 
\\ 
\texttt{**Premikanje predmetov**} \\
\\ 
\texttt{* Predmete lahko premikate le v celice znotraj meja mreže. Ciljna celica mora biti prazna, tj. vsebovati mora samo simbol '◌'.} \\
\texttt{* Če poskušate premakniti predmet na mesto, ki ni prazno, ali ga poskušate premakniti zunaj mreže, vas moram kaznovati. Dobili boste še en poskus.} \\
\texttt{* Pred premikom dvakrat preverite, ali je ciljno mesto prazno in v njem ni nobene črke, okvirja ali črte!} \\
\\ 
\\ 
\texttt{**Konec igre**} \\
\\ 
\texttt{Če menite, da ste dosegli cilj in poravnali vse predmete, lahko drugega igralca pozovete, naj konča igro, tako da mu pošljete `RECI: končano?`. Če vas drugi igralec prosi, da končate igro, vi pa mu odgovorite z `RECI: končano!`, se igra konča.} \\
\\ 
\texttt{Oba igralca zmagata, če se igra konča v 20 krogih, pri čemer je en krog definiran tako, da dva igralca pošljeta vsak po en veljaven ukaz.} \\
\\ 
\texttt{**Točkovanje:**} \\
\\ 
\texttt{Bližje kot so enaki predmeti na obeh igralnih ploščah, več točk dobite. Kazni zmanjšajo število točk. Ali lahko premagate rekord?} \\
\\ 
\texttt{Drugi igralec je začel igro s tem sporočilom:} \\
\\ 
\texttt{"Pozdravljeni! Dogovoriva se o naslednji ciljni konfiguraciji: W:(1,1), I:(2,1), T:(3,1), C:(4,1), H:(5,1). Ali se strinjaš?"} \\
\\ 
\texttt{Kakšen je vaš prvi ukaz?} \\
            }
        }
    }
    & \\ \\

    \theutterance \stepcounter{utterance}  
    & & & \multicolumn{4}{p{0.6\linewidth}}{
        \cellcolor[rgb]{0.9,0.9,0.9}{
            \makecell[{{p{\linewidth}}}]{
                \texttt{\tiny{[GM$\langle$P2]}}
                \texttt{RECI: Strinjam se s ciljno konfiguracijo: W:(1,1), I:(2,1), T:(3,1), C:(4,1), H:(5,1).} \\
            }
        }
    }
     \\ \\

    \theutterance \stepcounter{utterance}  
    & & & \multicolumn{2}{p{0.3\linewidth}}{
        \cellcolor[rgb]{0.9,0.9,0.9}{
            \makecell[{{p{\linewidth}}}]{
                \texttt{\tiny{[GM$|$GM]}}
                \texttt{RECI: Strinjam se s ciljno konfiguracijo: W:(1,1), I:(2,1), T:(3,1), C:(4,1), H:(5,1).} \\
            }
        }
    }
    & & \\ \\

    \theutterance \stepcounter{utterance}  
    & & \multicolumn{4}{p{0.6\linewidth}}{
        \cellcolor[rgb]{0.9,0.9,0.9}{
            \makecell[{{p{\linewidth}}}]{
                \texttt{\tiny{[P1$\langle$GM]}}
                \texttt{Vaše sporočilo je bilo posredovano drugemu igralcu.} \\
\\ 
\\ 
\texttt{Trenutno igrate krog 1 od največ 20.} \\
\texttt{Skupaj ste zbrali 1 od 12 kazni.} \\
\texttt{Drugi igralec je poslal to sporočilo:} \\
\\ 
\\ 
\texttt{"Strinjam se s ciljno konfiguracijo: W:(1,1), I:(2,1), T:(3,1), C:(4,1), H:(5,1)."} \\
\\ 
\\ 
\texttt{Kakšen je vaš naslednji ukaz?} \\
            }
        }
    }
    & \\ \\

    \theutterance \stepcounter{utterance}  
    & \multicolumn{4}{p{0.6\linewidth}}{
        \cellcolor[rgb]{0.9,0.9,0.9}{
            \makecell[{{p{\linewidth}}}]{
                \texttt{\tiny{[P1$\rangle$GM]}}
                \texttt{PREMAKNI SE: W, (1, 1)} \\
            }
        }
    }
    & & \\ \\

    \theutterance \stepcounter{utterance}  
    & & & \multicolumn{2}{p{0.3\linewidth}}{
        \cellcolor[rgb]{0.9,0.9,0.9}{
            \makecell[{{p{\linewidth}}}]{
                \texttt{\tiny{[GM$|$GM]}}
                \texttt{PREMAKNI SE: W, (1, 1)} \\
            }
        }
    }
    & & \\ \\

    \theutterance \stepcounter{utterance}  
    & & & \multicolumn{2}{p{0.3\linewidth}}{
        \cellcolor[rgb]{0.9,0.9,0.9}{
            \makecell[{{p{\linewidth}}}]{
                \texttt{\tiny{[GM$|$GM]}}
                \texttt{W moved to (1, 1)} \\
            }
        }
    }
    & & \\ \\

    \theutterance \stepcounter{utterance}  
    & & \multicolumn{4}{p{0.6\linewidth}}{
        \cellcolor[rgb]{0.9,0.9,0.9}{
            \makecell[{{p{\linewidth}}}]{
                \texttt{\tiny{[GM$\rangle$P2]}}
                \texttt{Vaše sporočilo je bilo posredovano drugemu igralcu.} \\
\\ 
\\ 
\texttt{Trenutno igrate krog 2 od največ 20.} \\
\texttt{Skupaj ste zbrali 1 od 12 kazni.} \\
\texttt{Drugi igralec je premaknil predmet na svoji mreži.} \\
\\ 
\\ 
\texttt{Kakšen je vaš naslednji ukaz?} \\
            }
        }
    }
    & \\ \\

    \theutterance \stepcounter{utterance}  
    & & & \multicolumn{4}{p{0.6\linewidth}}{
        \cellcolor[rgb]{0.9,0.9,0.9}{
            \makecell[{{p{\linewidth}}}]{
                \texttt{\tiny{[GM$\langle$P2]}}
                \texttt{PREMAKNI SE: W, (1, 1)} \\
            }
        }
    }
     \\ \\

    \theutterance \stepcounter{utterance}  
    & & & \multicolumn{2}{p{0.3\linewidth}}{
        \cellcolor[rgb]{0.9,0.9,0.9}{
            \makecell[{{p{\linewidth}}}]{
                \texttt{\tiny{[GM$|$GM]}}
                \texttt{PREMAKNI SE: W, (1, 1)} \\
            }
        }
    }
    & & \\ \\

    \theutterance \stepcounter{utterance}  
    & & & \multicolumn{2}{p{0.3\linewidth}}{
        \cellcolor[rgb]{0.9,0.9,0.9}{
            \makecell[{{p{\linewidth}}}]{
                \texttt{\tiny{[GM$|$GM]}}
                \texttt{W moved to (1, 1)} \\
            }
        }
    }
    & & \\ \\

    \theutterance \stepcounter{utterance}  
    & & \multicolumn{4}{p{0.6\linewidth}}{
        \cellcolor[rgb]{0.9,0.9,0.9}{
            \makecell[{{p{\linewidth}}}]{
                \texttt{\tiny{[P1$\langle$GM]}}
                \texttt{Uspešno premaknjen 'W' v (1,1). Vaša posodobljena mreža je videti takole:} \\
\texttt{```} \\
\texttt{ 1234567} \\
\texttt{╔═╤═╤═══╗} \\
\texttt{║W│◌│◌◌◌║ 1} \\
\texttt{╟─┴─┼─┬─╢ 2} \\
\texttt{║◌◌◌│◌│T║ 3} \\
\texttt{║◌H◌├─┘◌║ 4} \\
\texttt{║◌┌─┘◌◌◌║ 5} \\
\texttt{╟─┤◌◌◌◌◌║ 6} \\
\texttt{║I│◌C◌◌◌║ 7} \\
\texttt{╚═╧═════╝} \\
\texttt{```} \\
\texttt{Trenutno igrate krog 2 od največ 20.} \\
\texttt{Skupaj ste zbrali 1 od 12 kazni.} \\
\texttt{Drugi igralec je premaknil predmet na svoji mreži.} \\
\\ 
\\ 
\texttt{Kakšen je vaš naslednji ukaz?} \\
            }
        }
    }
    & \\ \\

    \theutterance \stepcounter{utterance}  
    & \multicolumn{4}{p{0.6\linewidth}}{
        \cellcolor[rgb]{0.9,0.9,0.9}{
            \makecell[{{p{\linewidth}}}]{
                \texttt{\tiny{[P1$\rangle$GM]}}
                \texttt{PREMAKNI SE: I, (2, 1)} \\
            }
        }
    }
    & & \\ \\

    \theutterance \stepcounter{utterance}  
    & & & \multicolumn{2}{p{0.3\linewidth}}{
        \cellcolor[rgb]{0.9,0.9,0.9}{
            \makecell[{{p{\linewidth}}}]{
                \texttt{\tiny{[GM$|$GM]}}
                \texttt{PREMAKNI SE: I, (2, 1)} \\
            }
        }
    }
    & & \\ \\

    \theutterance \stepcounter{utterance}  
    & & & \multicolumn{2}{p{0.3\linewidth}}{
        \cellcolor[rgb]{0.9,0.9,0.9}{
            \makecell[{{p{\linewidth}}}]{
                \texttt{\tiny{[GM$|$GM]}}
                \texttt{Kazen: (2,1) ni prazen, vendar vsebuje '│'.} \\
\\ 
\texttt{Skupaj ste zbrali 2 od 12 kazni.} \\
\texttt{Prosim, poskusite znova!} \\
            }
        }
    }
    & & \\ \\

    \theutterance \stepcounter{utterance}  
    & & \multicolumn{4}{p{0.6\linewidth}}{
        \cellcolor[rgb]{0.9,0.9,0.9}{
            \makecell[{{p{\linewidth}}}]{
                \texttt{\tiny{[P1$\langle$GM]}}
                \texttt{Kazen: (2,1) ni prazen, vendar vsebuje '│'.} \\
\\ 
\texttt{Skupaj ste zbrali 2 od 12 kazni.} \\
\texttt{Prosim, poskusite znova!} \\
            }
        }
    }
    & \\ \\

    \theutterance \stepcounter{utterance}  
    & \multicolumn{4}{p{0.6\linewidth}}{
        \cellcolor[rgb]{0.9,0.9,0.9}{
            \makecell[{{p{\linewidth}}}]{
                \texttt{\tiny{[P1$\rangle$GM]}}
                \texttt{PREMAKNI SE: I, (2, 1)} \\
\\ 
\texttt{(Opomba: Napaka je bila v tem, da sem v prejšnjem koraku neposredno poskušal premakniti predmet na koordinato, ki je bila v mojem razumevanju mreže vključena v okvir/črte. Ker sem dobil kazen, poskušam znova, vendar moram biti previden pri branju koordinat. Na moji mreži je (2,1) celica med stolpci 1 in 2 na vrstici 1. Če je (2,1) v celoti znotraj mreže in brez črt, bo uspešno. Če pa je (2,1) dejansko črta, moram izbrati drugo koordinato. Ker pa sem ukaz že poslal, čakamo na rezultat.)} \\
\\ 
\texttt{*Ureditev: Ker sem že poslal ukaz, ne morem dodati opombe. Čakam na odziv vodje igre.*} \\
            }
        }
    }
    & & \\ \\

    \theutterance \stepcounter{utterance}  
    & & & \multicolumn{2}{p{0.3\linewidth}}{
        \cellcolor[rgb]{0.9,0.9,0.9}{
            \makecell[{{p{\linewidth}}}]{
                \texttt{\tiny{[GM$|$GM]}}
                \texttt{PREMAKNI SE: I, (2, 1)} \\
\\ 
\texttt{(Opomba: Napaka je bila v tem, da sem v prejšnjem koraku neposredno poskušal premakniti predmet na koordinato, ki je bila v mojem razumevanju mreže vključena v okvir/črte. Ker sem dobil kazen, poskušam znova, vendar moram biti previden pri branju koordinat. Na moji mreži je (2,1) celica med stolpci 1 in 2 na vrstici 1. Če je (2,1) v celoti znotraj mreže in brez črt, bo uspešno. Če pa je (2,1) dejansko črta, moram izbrati drugo koordinato. Ker pa sem ukaz že poslal, čakamo na rezultat.)} \\
\\ 
\texttt{*Ureditev: Ker sem že poslal ukaz, ne morem dodati opombe. Čakam na odziv vodje igre.*} \\
            }
        }
    }
    & & \\ \\

    \theutterance \stepcounter{utterance}  
    & & & \multicolumn{2}{p{0.3\linewidth}}{
        \cellcolor[rgb]{0.9,0.9,0.9}{
            \makecell[{{p{\linewidth}}}]{
                \texttt{\tiny{[GM$|$GM]}}
                \texttt{Invalid format: tail is not empty: '} \\
\\ 
\texttt{(Opomba: Napaka je bila v tem, da sem v prejšnjem koraku neposredno poskušal premakniti predmet na koordinato, ki je bila v mojem razumevanju mreže vključena v okvir/črte. Ker sem dobil kazen, poskušam znova, vendar moram biti previden pri branju koordinat. Na moji mreži je (2,1) celica med stolpci 1 in 2 na vrstici 1. Če je (2,1) v celoti znotraj mreže in brez črt, bo uspešno. Če pa je (2,1) dejansko črta, moram izbrati drugo koordinato. Ker pa sem ukaz že poslal, čakamo na rezultat.)} \\
\\ 
\texttt{*Ureditev: Ker sem že poslal ukaz, ne morem dodati opombe. Čakam na odziv vodje igre.*} \\
            }
        }
    }
    & & \\ \\

    \theutterance \stepcounter{utterance}  
    & & \multicolumn{4}{p{0.6\linewidth}}{
        \cellcolor[rgb]{0.9,0.9,0.9}{
            \makecell[{{p{\linewidth}}}]{
                \texttt{\tiny{[P1$\langle$GM]}}
                \texttt{Kazen: Vaše sporočilo ne sme vsebovati ničesar po ukazu!} \\
\\ 
\texttt{Prepričajte se, da vaš odgovor vsebuje samo `RECI: <SPOROČILO>` ali `PREMAKNI SE: <PREDMET>, (<X>, <Y>)` in nič drugega!} \\
\\ 
\texttt{Skupaj ste zbrali 3 od 12 kazni.} \\
\texttt{Prosim, poskusite znova!} \\
            }
        }
    }
    & \\ \\

    \theutterance \stepcounter{utterance}  
    & \multicolumn{4}{p{0.6\linewidth}}{
        \cellcolor[rgb]{0.9,0.9,0.9}{
            \makecell[{{p{\linewidth}}}]{
                \texttt{\tiny{[P1$\rangle$GM]}}
                \texttt{PREMAKNI SE: I, (1, 2)} \\
            }
        }
    }
    & & \\ \\

    \theutterance \stepcounter{utterance}  
    & & & \multicolumn{2}{p{0.3\linewidth}}{
        \cellcolor[rgb]{0.9,0.9,0.9}{
            \makecell[{{p{\linewidth}}}]{
                \texttt{\tiny{[GM$|$GM]}}
                \texttt{PREMAKNI SE: I, (1, 2)} \\
            }
        }
    }
    & & \\ \\

    \theutterance \stepcounter{utterance}  
    & & & \multicolumn{2}{p{0.3\linewidth}}{
        \cellcolor[rgb]{0.9,0.9,0.9}{
            \makecell[{{p{\linewidth}}}]{
                \texttt{\tiny{[GM$|$GM]}}
                \texttt{Kazen: (1,2) ni prazen, vendar vsebuje '─'.} \\
\\ 
\texttt{Skupaj ste zbrali 4 od 12 kazni.} \\
\texttt{Prosim, poskusite znova!} \\
            }
        }
    }
    & & \\ \\

    \theutterance \stepcounter{utterance}  
    & & \multicolumn{4}{p{0.6\linewidth}}{
        \cellcolor[rgb]{0.9,0.9,0.9}{
            \makecell[{{p{\linewidth}}}]{
                \texttt{\tiny{[P1$\langle$GM]}}
                \texttt{Kazen: (1,2) ni prazen, vendar vsebuje '─'.} \\
\\ 
\texttt{Skupaj ste zbrali 4 od 12 kazni.} \\
\texttt{Prosim, poskusite znova!} \\
            }
        }
    }
    & \\ \\

    \theutterance \stepcounter{utterance}  
    & \multicolumn{4}{p{0.6\linewidth}}{
        \cellcolor[rgb]{0.9,0.9,0.9}{
            \makecell[{{p{\linewidth}}}]{
                \texttt{\tiny{[P1$\rangle$GM]}}
                \texttt{PREMAKNI SE: I, (1, 3)} \\
            }
        }
    }
    & & \\ \\

    \theutterance \stepcounter{utterance}  
    & & & \multicolumn{2}{p{0.3\linewidth}}{
        \cellcolor[rgb]{0.9,0.9,0.9}{
            \makecell[{{p{\linewidth}}}]{
                \texttt{\tiny{[GM$|$GM]}}
                \texttt{PREMAKNI SE: I, (1, 3)} \\
            }
        }
    }
    & & \\ \\

    \theutterance \stepcounter{utterance}  
    & & & \multicolumn{2}{p{0.3\linewidth}}{
        \cellcolor[rgb]{0.9,0.9,0.9}{
            \makecell[{{p{\linewidth}}}]{
                \texttt{\tiny{[GM$|$GM]}}
                \texttt{I moved to (1, 3)} \\
            }
        }
    }
    & & \\ \\

    \theutterance \stepcounter{utterance}  
    & & \multicolumn{4}{p{0.6\linewidth}}{
        \cellcolor[rgb]{0.9,0.9,0.9}{
            \makecell[{{p{\linewidth}}}]{
                \texttt{\tiny{[GM$\rangle$P2]}}
                \texttt{Uspešno premaknjen 'W' v (1,1). Vaša posodobljena mreža je videti takole:} \\
\texttt{```} \\
\texttt{ 1234567} \\
\texttt{╔═╤═╤═══╗} \\
\texttt{║W│◌│◌◌◌║ 1} \\
\texttt{╟─┴─┼─┬─╢ 2} \\
\texttt{║◌◌◌│◌│◌║ 3} \\
\texttt{║◌◌C├─┘◌║ 4} \\
\texttt{║◌┌─┘◌T◌║ 5} \\
\texttt{╟─┤◌◌◌◌H║ 6} \\
\texttt{║◌│◌I◌◌◌║ 7} \\
\texttt{╚═╧═════╝} \\
\texttt{```} \\
\texttt{Trenutno igrate krog 3 od največ 20.} \\
\texttt{Skupaj ste zbrali 4 od 12 kazni.} \\
\texttt{Drugi igralec je premaknil predmet na svoji mreži.} \\
\\ 
\\ 
\texttt{Kakšen je vaš naslednji ukaz?} \\
            }
        }
    }
    & \\ \\

    \theutterance \stepcounter{utterance}  
    & & & \multicolumn{4}{p{0.6\linewidth}}{
        \cellcolor[rgb]{0.9,0.9,0.9}{
            \makecell[{{p{\linewidth}}}]{
                \texttt{\tiny{[GM$\langle$P2]}}
                \texttt{PREMAKNI SE: I, (2, 1)} \\
            }
        }
    }
     \\ \\

    \theutterance \stepcounter{utterance}  
    & & & \multicolumn{2}{p{0.3\linewidth}}{
        \cellcolor[rgb]{0.9,0.9,0.9}{
            \makecell[{{p{\linewidth}}}]{
                \texttt{\tiny{[GM$|$GM]}}
                \texttt{PREMAKNI SE: I, (2, 1)} \\
            }
        }
    }
    & & \\ \\

    \theutterance \stepcounter{utterance}  
    & & & \multicolumn{2}{p{0.3\linewidth}}{
        \cellcolor[rgb]{0.9,0.9,0.9}{
            \makecell[{{p{\linewidth}}}]{
                \texttt{\tiny{[GM$|$GM]}}
                \texttt{Kazen: (2,1) ni prazen, vendar vsebuje '│'.} \\
\\ 
\texttt{Skupaj ste zbrali 5 od 12 kazni.} \\
\texttt{Prosim, poskusite znova!} \\
            }
        }
    }
    & & \\ \\

    \theutterance \stepcounter{utterance}  
    & & \multicolumn{4}{p{0.6\linewidth}}{
        \cellcolor[rgb]{0.9,0.9,0.9}{
            \makecell[{{p{\linewidth}}}]{
                \texttt{\tiny{[GM$\rangle$P2]}}
                \texttt{Kazen: (2,1) ni prazen, vendar vsebuje '│'.} \\
\\ 
\texttt{Skupaj ste zbrali 5 od 12 kazni.} \\
\texttt{Prosim, poskusite znova!} \\
            }
        }
    }
    & \\ \\

    \theutterance \stepcounter{utterance}  
    & & & \multicolumn{4}{p{0.6\linewidth}}{
        \cellcolor[rgb]{0.9,0.9,0.9}{
            \makecell[{{p{\linewidth}}}]{
                \texttt{\tiny{[GM$\langle$P2]}}
                \texttt{PREMAKNI SE: I, (2, 7)} \\
            }
        }
    }
     \\ \\

    \theutterance \stepcounter{utterance}  
    & & & \multicolumn{2}{p{0.3\linewidth}}{
        \cellcolor[rgb]{0.9,0.9,0.9}{
            \makecell[{{p{\linewidth}}}]{
                \texttt{\tiny{[GM$|$GM]}}
                \texttt{PREMAKNI SE: I, (2, 7)} \\
            }
        }
    }
    & & \\ \\

    \theutterance \stepcounter{utterance}  
    & & & \multicolumn{2}{p{0.3\linewidth}}{
        \cellcolor[rgb]{0.9,0.9,0.9}{
            \makecell[{{p{\linewidth}}}]{
                \texttt{\tiny{[GM$|$GM]}}
                \texttt{Kazen: (2,7) ni prazen, vendar vsebuje '│'.} \\
\\ 
\texttt{Skupaj ste zbrali 6 od 12 kazni.} \\
\texttt{Prosim, poskusite znova!} \\
            }
        }
    }
    & & \\ \\

    \theutterance \stepcounter{utterance}  
    & & \multicolumn{4}{p{0.6\linewidth}}{
        \cellcolor[rgb]{0.9,0.9,0.9}{
            \makecell[{{p{\linewidth}}}]{
                \texttt{\tiny{[GM$\rangle$P2]}}
                \texttt{Kazen: (2,7) ni prazen, vendar vsebuje '│'.} \\
\\ 
\texttt{Skupaj ste zbrali 6 od 12 kazni.} \\
\texttt{Prosim, poskusite znova!} \\
            }
        }
    }
    & \\ \\

    \theutterance \stepcounter{utterance}  
    & & & \multicolumn{4}{p{0.6\linewidth}}{
        \cellcolor[rgb]{0.9,0.9,0.9}{
            \makecell[{{p{\linewidth}}}]{
                \texttt{\tiny{[GM$\langle$P2]}}
                \texttt{PREMAKNI SE: I, (3, 7)} \\
            }
        }
    }
     \\ \\

    \theutterance \stepcounter{utterance}  
    & & & \multicolumn{2}{p{0.3\linewidth}}{
        \cellcolor[rgb]{0.9,0.9,0.9}{
            \makecell[{{p{\linewidth}}}]{
                \texttt{\tiny{[GM$|$GM]}}
                \texttt{PREMAKNI SE: I, (3, 7)} \\
            }
        }
    }
    & & \\ \\

    \theutterance \stepcounter{utterance}  
    & & & \multicolumn{2}{p{0.3\linewidth}}{
        \cellcolor[rgb]{0.9,0.9,0.9}{
            \makecell[{{p{\linewidth}}}]{
                \texttt{\tiny{[GM$|$GM]}}
                \texttt{I moved to (3, 7)} \\
            }
        }
    }
    & & \\ \\

    \theutterance \stepcounter{utterance}  
    & & \multicolumn{4}{p{0.6\linewidth}}{
        \cellcolor[rgb]{0.9,0.9,0.9}{
            \makecell[{{p{\linewidth}}}]{
                \texttt{\tiny{[P1$\langle$GM]}}
                \texttt{Uspešno premaknjen 'I' v (1,3). Vaša posodobljena mreža je videti takole:} \\
\texttt{```} \\
\texttt{ 1234567} \\
\texttt{╔═╤═╤═══╗} \\
\texttt{║W│◌│◌◌◌║ 1} \\
\texttt{╟─┴─┼─┬─╢ 2} \\
\texttt{║I◌◌│◌│T║ 3} \\
\texttt{║◌H◌├─┘◌║ 4} \\
\texttt{║◌┌─┘◌◌◌║ 5} \\
\texttt{╟─┤◌◌◌◌◌║ 6} \\
\texttt{║◌│◌C◌◌◌║ 7} \\
\texttt{╚═╧═════╝} \\
\texttt{```} \\
\texttt{Trenutno igrate krog 3 od največ 20.} \\
\texttt{Skupaj ste zbrali 6 od 12 kazni.} \\
\texttt{Drugi igralec je premaknil predmet na svoji mreži.} \\
\\ 
\\ 
\texttt{Kakšen je vaš naslednji ukaz?} \\
            }
        }
    }
    & \\ \\

    \theutterance \stepcounter{utterance}  
    & \multicolumn{4}{p{0.6\linewidth}}{
        \cellcolor[rgb]{0.9,0.9,0.9}{
            \makecell[{{p{\linewidth}}}]{
                \texttt{\tiny{[P1$\rangle$GM]}}
                \texttt{PREMAKNI SE: H, (2, 4)} \\
            }
        }
    }
    & & \\ \\

    \theutterance \stepcounter{utterance}  
    & & & \multicolumn{2}{p{0.3\linewidth}}{
        \cellcolor[rgb]{0.9,0.9,0.9}{
            \makecell[{{p{\linewidth}}}]{
                \texttt{\tiny{[GM$|$GM]}}
                \texttt{PREMAKNI SE: H, (2, 4)} \\
            }
        }
    }
    & & \\ \\

    \theutterance \stepcounter{utterance}  
    & & & \multicolumn{2}{p{0.3\linewidth}}{
        \cellcolor[rgb]{0.9,0.9,0.9}{
            \makecell[{{p{\linewidth}}}]{
                \texttt{\tiny{[GM$|$GM]}}
                \texttt{Kazen: (2,4) ni prazen, vendar vsebuje 'H'.} \\
\\ 
\texttt{Skupaj ste zbrali 7 od 12 kazni.} \\
\texttt{Prosim, poskusite znova!} \\
            }
        }
    }
    & & \\ \\

    \theutterance \stepcounter{utterance}  
    & & \multicolumn{4}{p{0.6\linewidth}}{
        \cellcolor[rgb]{0.9,0.9,0.9}{
            \makecell[{{p{\linewidth}}}]{
                \texttt{\tiny{[P1$\langle$GM]}}
                \texttt{Kazen: (2,4) ni prazen, vendar vsebuje 'H'.} \\
\\ 
\texttt{Skupaj ste zbrali 7 od 12 kazni.} \\
\texttt{Prosim, poskusite znova!} \\
            }
        }
    }
    & \\ \\

    \theutterance \stepcounter{utterance}  
    & \multicolumn{4}{p{0.6\linewidth}}{
        \cellcolor[rgb]{0.9,0.9,0.9}{
            \makecell[{{p{\linewidth}}}]{
                \texttt{\tiny{[P1$\rangle$GM]}}
                \texttt{PREMAKNI SE: H, (2, 4)} \\
\\ 
\texttt{(Opomba: Napaka je bila v tem, da sem poskušal premakniti predmet na njegovo trenutno lokacijo. Ker sem dobil kazen, moram izbrati novo lokacijo. Ker je H že na (2,4), moram izbrati prazno celico. Na moji mreži je (2,4) trenutno H. Poskusim premakniti H na (2,3).)} \\
\\ 
\texttt{PREMAKNI SE: H, (2, 3)} \\
            }
        }
    }
    & & \\ \\

    \theutterance \stepcounter{utterance}  
    & & & \multicolumn{2}{p{0.3\linewidth}}{
        \cellcolor[rgb]{0.9,0.9,0.9}{
            \makecell[{{p{\linewidth}}}]{
                \texttt{\tiny{[GM$|$GM]}}
                \texttt{PREMAKNI SE: H, (2, 4)} \\
\\ 
\texttt{(Opomba: Napaka je bila v tem, da sem poskušal premakniti predmet na njegovo trenutno lokacijo. Ker sem dobil kazen, moram izbrati novo lokacijo. Ker je H že na (2,4), moram izbrati prazno celico. Na moji mreži je (2,4) trenutno H. Poskusim premakniti H na (2,3).)} \\
\\ 
\texttt{PREMAKNI SE: H, (2, 3)} \\
            }
        }
    }
    & & \\ \\

    \theutterance \stepcounter{utterance}  
    & & & \multicolumn{2}{p{0.3\linewidth}}{
        \cellcolor[rgb]{0.9,0.9,0.9}{
            \makecell[{{p{\linewidth}}}]{
                \texttt{\tiny{[GM$|$GM]}}
                \texttt{Invalid format: head is not empty: 'PREMAKNI SE: H, (2, 4)} \\
\\ 
\texttt{(Opomba: Napaka je bila v tem, da sem poskušal premakniti predmet na njegovo trenutno lokacijo. Ker sem dobil kazen, moram izbrati novo lokacijo. Ker je H že na (2,4), moram izbrati prazno celico. Na moji mreži je (2,4) trenutno H. Poskusim premakniti H na (2,3).)} \\
\\ 
\texttt{'} \\
            }
        }
    }
    & & \\ \\

    \theutterance \stepcounter{utterance}  
    & & \multicolumn{4}{p{0.6\linewidth}}{
        \cellcolor[rgb]{0.9,0.9,0.9}{
            \makecell[{{p{\linewidth}}}]{
                \texttt{\tiny{[P1$\langle$GM]}}
                \texttt{Kazen: Vaše sporočilo ne sme vsebovati ničesar pred ukazom!} \\
\\ 
\texttt{Prepričajte se, da vaš odgovor vsebuje samo `RECI: <SPOROČILO>` ali `PREMAKNI SE: <PREDMET>, (<X>, <Y>)` in nič drugega!} \\
\\ 
\texttt{Skupaj ste zbrali 8 od 12 kazni.} \\
\texttt{Prosim, poskusite znova!} \\
            }
        }
    }
    & \\ \\

    \theutterance \stepcounter{utterance}  
    & \multicolumn{4}{p{0.6\linewidth}}{
        \cellcolor[rgb]{0.9,0.9,0.9}{
            \makecell[{{p{\linewidth}}}]{
                \texttt{\tiny{[P1$\rangle$GM]}}
                \texttt{PREMAKNI SE: H, (2, 3)} \\
            }
        }
    }
    & & \\ \\

    \theutterance \stepcounter{utterance}  
    & & & \multicolumn{2}{p{0.3\linewidth}}{
        \cellcolor[rgb]{0.9,0.9,0.9}{
            \makecell[{{p{\linewidth}}}]{
                \texttt{\tiny{[GM$|$GM]}}
                \texttt{PREMAKNI SE: H, (2, 3)} \\
            }
        }
    }
    & & \\ \\

    \theutterance \stepcounter{utterance}  
    & & & \multicolumn{2}{p{0.3\linewidth}}{
        \cellcolor[rgb]{0.9,0.9,0.9}{
            \makecell[{{p{\linewidth}}}]{
                \texttt{\tiny{[GM$|$GM]}}
                \texttt{H moved to (2, 3)} \\
            }
        }
    }
    & & \\ \\

    \theutterance \stepcounter{utterance}  
    & & \multicolumn{4}{p{0.6\linewidth}}{
        \cellcolor[rgb]{0.9,0.9,0.9}{
            \makecell[{{p{\linewidth}}}]{
                \texttt{\tiny{[GM$\rangle$P2]}}
                \texttt{Uspešno premaknjen 'I' v (3,7). Vaša posodobljena mreža je videti takole:} \\
\texttt{```} \\
\texttt{ 1234567} \\
\texttt{╔═╤═╤═══╗} \\
\texttt{║W│◌│◌◌◌║ 1} \\
\texttt{╟─┴─┼─┬─╢ 2} \\
\texttt{║◌◌◌│◌│◌║ 3} \\
\texttt{║◌◌C├─┘◌║ 4} \\
\texttt{║◌┌─┘◌T◌║ 5} \\
\texttt{╟─┤◌◌◌◌H║ 6} \\
\texttt{║◌│I◌◌◌◌║ 7} \\
\texttt{╚═╧═════╝} \\
\texttt{```} \\
\texttt{Trenutno igrate krog 4 od največ 20.} \\
\texttt{Skupaj ste zbrali 8 od 12 kazni.} \\
\texttt{Drugi igralec je premaknil predmet na svoji mreži.} \\
\\ 
\\ 
\texttt{Kakšen je vaš naslednji ukaz?} \\
            }
        }
    }
    & \\ \\

    \theutterance \stepcounter{utterance}  
    & & & \multicolumn{4}{p{0.6\linewidth}}{
        \cellcolor[rgb]{0.9,0.9,0.9}{
            \makecell[{{p{\linewidth}}}]{
                \texttt{\tiny{[GM$\langle$P2]}}
                \texttt{PREMAKNI SE: T, (3, 5)} \\
            }
        }
    }
     \\ \\

    \theutterance \stepcounter{utterance}  
    & & & \multicolumn{2}{p{0.3\linewidth}}{
        \cellcolor[rgb]{0.9,0.9,0.9}{
            \makecell[{{p{\linewidth}}}]{
                \texttt{\tiny{[GM$|$GM]}}
                \texttt{PREMAKNI SE: T, (3, 5)} \\
            }
        }
    }
    & & \\ \\

    \theutterance \stepcounter{utterance}  
    & & & \multicolumn{2}{p{0.3\linewidth}}{
        \cellcolor[rgb]{0.9,0.9,0.9}{
            \makecell[{{p{\linewidth}}}]{
                \texttt{\tiny{[GM$|$GM]}}
                \texttt{Kazen: (3,5) ni prazen, vendar vsebuje '─'.} \\
\\ 
\texttt{Skupaj ste zbrali 9 od 12 kazni.} \\
\texttt{Prosim, poskusite znova!} \\
            }
        }
    }
    & & \\ \\

    \theutterance \stepcounter{utterance}  
    & & \multicolumn{4}{p{0.6\linewidth}}{
        \cellcolor[rgb]{0.9,0.9,0.9}{
            \makecell[{{p{\linewidth}}}]{
                \texttt{\tiny{[GM$\rangle$P2]}}
                \texttt{Kazen: (3,5) ni prazen, vendar vsebuje '─'.} \\
\\ 
\texttt{Skupaj ste zbrali 9 od 12 kazni.} \\
\texttt{Prosim, poskusite znova!} \\
            }
        }
    }
    & \\ \\

    \theutterance \stepcounter{utterance}  
    & & & \multicolumn{4}{p{0.6\linewidth}}{
        \cellcolor[rgb]{0.9,0.9,0.9}{
            \makecell[{{p{\linewidth}}}]{
                \texttt{\tiny{[GM$\langle$P2]}}
                \texttt{PREMAKNI SE: T, (5, 5)} \\
            }
        }
    }
     \\ \\

    \theutterance \stepcounter{utterance}  
    & & & \multicolumn{2}{p{0.3\linewidth}}{
        \cellcolor[rgb]{0.9,0.9,0.9}{
            \makecell[{{p{\linewidth}}}]{
                \texttt{\tiny{[GM$|$GM]}}
                \texttt{PREMAKNI SE: T, (5, 5)} \\
            }
        }
    }
    & & \\ \\

    \theutterance \stepcounter{utterance}  
    & & & \multicolumn{2}{p{0.3\linewidth}}{
        \cellcolor[rgb]{0.9,0.9,0.9}{
            \makecell[{{p{\linewidth}}}]{
                \texttt{\tiny{[GM$|$GM]}}
                \texttt{T moved to (5, 5)} \\
            }
        }
    }
    & & \\ \\

    \theutterance \stepcounter{utterance}  
    & & \multicolumn{4}{p{0.6\linewidth}}{
        \cellcolor[rgb]{0.9,0.9,0.9}{
            \makecell[{{p{\linewidth}}}]{
                \texttt{\tiny{[P1$\langle$GM]}}
                \texttt{Uspešno premaknjen 'H' v (2,3). Vaša posodobljena mreža je videti takole:} \\
\texttt{```} \\
\texttt{ 1234567} \\
\texttt{╔═╤═╤═══╗} \\
\texttt{║W│◌│◌◌◌║ 1} \\
\texttt{╟─┴─┼─┬─╢ 2} \\
\texttt{║IH◌│◌│T║ 3} \\
\texttt{║◌◌◌├─┘◌║ 4} \\
\texttt{║◌┌─┘◌◌◌║ 5} \\
\texttt{╟─┤◌◌◌◌◌║ 6} \\
\texttt{║◌│◌C◌◌◌║ 7} \\
\texttt{╚═╧═════╝} \\
\texttt{```} \\
\texttt{Trenutno igrate krog 4 od največ 20.} \\
\texttt{Skupaj ste zbrali 9 od 12 kazni.} \\
\texttt{Drugi igralec je premaknil predmet na svoji mreži.} \\
\\ 
\\ 
\texttt{Kakšen je vaš naslednji ukaz?} \\
            }
        }
    }
    & \\ \\

    \theutterance \stepcounter{utterance}  
    & \multicolumn{4}{p{0.6\linewidth}}{
        \cellcolor[rgb]{0.9,0.9,0.9}{
            \makecell[{{p{\linewidth}}}]{
                \texttt{\tiny{[P1$\rangle$GM]}}
                \texttt{PREMAKNI SE: H, (2, 1)} \\
            }
        }
    }
    & & \\ \\

    \theutterance \stepcounter{utterance}  
    & & & \multicolumn{2}{p{0.3\linewidth}}{
        \cellcolor[rgb]{0.9,0.9,0.9}{
            \makecell[{{p{\linewidth}}}]{
                \texttt{\tiny{[GM$|$GM]}}
                \texttt{PREMAKNI SE: H, (2, 1)} \\
            }
        }
    }
    & & \\ \\

    \theutterance \stepcounter{utterance}  
    & & & \multicolumn{2}{p{0.3\linewidth}}{
        \cellcolor[rgb]{0.9,0.9,0.9}{
            \makecell[{{p{\linewidth}}}]{
                \texttt{\tiny{[GM$|$GM]}}
                \texttt{Kazen: (2,1) ni prazen, vendar vsebuje '│'.} \\
\\ 
\texttt{Skupaj ste zbrali 10 od 12 kazni.} \\
\texttt{Prosim, poskusite znova!} \\
            }
        }
    }
    & & \\ \\

    \theutterance \stepcounter{utterance}  
    & & \multicolumn{4}{p{0.6\linewidth}}{
        \cellcolor[rgb]{0.9,0.9,0.9}{
            \makecell[{{p{\linewidth}}}]{
                \texttt{\tiny{[P1$\langle$GM]}}
                \texttt{Kazen: (2,1) ni prazen, vendar vsebuje '│'.} \\
\\ 
\texttt{Skupaj ste zbrali 10 od 12 kazni.} \\
\texttt{Prosim, poskusite znova!} \\
            }
        }
    }
    & \\ \\

    \theutterance \stepcounter{utterance}  
    & \multicolumn{4}{p{0.6\linewidth}}{
        \cellcolor[rgb]{0.9,0.9,0.9}{
            \makecell[{{p{\linewidth}}}]{
                \texttt{\tiny{[P1$\rangle$GM]}}
                \texttt{PREMAKNI SE: H, (2, 2)} \\
\\ 
\texttt{(Opomba: Preverjam mrežo. Na vrstici 2 so črte ╟─┴─┼─┬─╢. Koordinata (2,2) je verjetno del črte. Moram najti prazno celico. Prazne celice so npr. (2,1) {-} ne, tam je │; (3,1) {-} da; (4,1) {-} da; (5,1) {-} da. Poskusim s (2,1) je bila napaka, poskusim s (3,1).)} \\
\\ 
\texttt{PREMAKNI SE: H, (3, 1)} \\
            }
        }
    }
    & & \\ \\

    \theutterance \stepcounter{utterance}  
    & & & \multicolumn{2}{p{0.3\linewidth}}{
        \cellcolor[rgb]{0.9,0.9,0.9}{
            \makecell[{{p{\linewidth}}}]{
                \texttt{\tiny{[GM$|$GM]}}
                \texttt{PREMAKNI SE: H, (2, 2)} \\
\\ 
\texttt{(Opomba: Preverjam mrežo. Na vrstici 2 so črte ╟─┴─┼─┬─╢. Koordinata (2,2) je verjetno del črte. Moram najti prazno celico. Prazne celice so npr. (2,1) {-} ne, tam je │; (3,1) {-} da; (4,1) {-} da; (5,1) {-} da. Poskusim s (2,1) je bila napaka, poskusim s (3,1).)} \\
\\ 
\texttt{PREMAKNI SE: H, (3, 1)} \\
            }
        }
    }
    & & \\ \\

    \theutterance \stepcounter{utterance}  
    & & & \multicolumn{2}{p{0.3\linewidth}}{
        \cellcolor[rgb]{0.9,0.9,0.9}{
            \makecell[{{p{\linewidth}}}]{
                \texttt{\tiny{[GM$|$GM]}}
                \texttt{Invalid format: head is not empty: 'PREMAKNI SE: H, (2, 2)} \\
\\ 
\texttt{(Opomba: Preverjam mrežo. Na vrstici 2 so črte ╟─┴─┼─┬─╢. Koordinata (2,2) je verjetno del črte. Moram najti prazno celico. Prazne celice so npr. (2,1) {-} ne, tam je │; (3,1) {-} da; (4,1) {-} da; (5,1) {-} da. Poskusim s (2,1) je bila napaka, poskusim s (3,1).)} \\
\\ 
\texttt{'} \\
            }
        }
    }
    & & \\ \\

    \theutterance \stepcounter{utterance}  
    & & \multicolumn{4}{p{0.6\linewidth}}{
        \cellcolor[rgb]{0.9,0.9,0.9}{
            \makecell[{{p{\linewidth}}}]{
                \texttt{\tiny{[P1$\langle$GM]}}
                \texttt{Kazen: Vaše sporočilo ne sme vsebovati ničesar pred ukazom!} \\
\\ 
\texttt{Prepričajte se, da vaš odgovor vsebuje samo `RECI: <SPOROČILO>` ali `PREMAKNI SE: <PREDMET>, (<X>, <Y>)` in nič drugega!} \\
\\ 
\texttt{Skupaj ste zbrali 11 od 12 kazni.} \\
\texttt{Prosim, poskusite znova!} \\
            }
        }
    }
    & \\ \\

    \theutterance \stepcounter{utterance}  
    & \multicolumn{4}{p{0.6\linewidth}}{
        \cellcolor[rgb]{0.9,0.9,0.9}{
            \makecell[{{p{\linewidth}}}]{
                \texttt{\tiny{[P1$\rangle$GM]}}
                \texttt{PREMAKNI SE: H, (3, 1)} \\
            }
        }
    }
    & & \\ \\

    \theutterance \stepcounter{utterance}  
    & & & \multicolumn{2}{p{0.3\linewidth}}{
        \cellcolor[rgb]{0.9,0.9,0.9}{
            \makecell[{{p{\linewidth}}}]{
                \texttt{\tiny{[GM$|$GM]}}
                \texttt{PREMAKNI SE: H, (3, 1)} \\
            }
        }
    }
    & & \\ \\

    \theutterance \stepcounter{utterance}  
    & & & \multicolumn{2}{p{0.3\linewidth}}{
        \cellcolor[rgb]{0.9,0.9,0.9}{
            \makecell[{{p{\linewidth}}}]{
                \texttt{\tiny{[GM$|$GM]}}
                \texttt{H moved to (3, 1)} \\
            }
        }
    }
    & & \\ \\

    \theutterance \stepcounter{utterance}  
    & & \multicolumn{4}{p{0.6\linewidth}}{
        \cellcolor[rgb]{0.9,0.9,0.9}{
            \makecell[{{p{\linewidth}}}]{
                \texttt{\tiny{[GM$\rangle$P2]}}
                \texttt{Uspešno premaknjen 'T' v (5,5). Vaša posodobljena mreža je videti takole:} \\
\texttt{```} \\
\texttt{ 1234567} \\
\texttt{╔═╤═╤═══╗} \\
\texttt{║W│◌│◌◌◌║ 1} \\
\texttt{╟─┴─┼─┬─╢ 2} \\
\texttt{║◌◌◌│◌│◌║ 3} \\
\texttt{║◌◌C├─┘◌║ 4} \\
\texttt{║◌┌─┘T◌◌║ 5} \\
\texttt{╟─┤◌◌◌◌H║ 6} \\
\texttt{║◌│I◌◌◌◌║ 7} \\
\texttt{╚═╧═════╝} \\
\texttt{```} \\
\texttt{Trenutno igrate krog 5 od največ 20.} \\
\texttt{Skupaj ste zbrali 11 od 12 kazni.} \\
\texttt{Drugi igralec je premaknil predmet na svoji mreži.} \\
\\ 
\\ 
\texttt{Kakšen je vaš naslednji ukaz?} \\
            }
        }
    }
    & \\ \\

    \theutterance \stepcounter{utterance}  
    & & & \multicolumn{4}{p{0.6\linewidth}}{
        \cellcolor[rgb]{0.9,0.9,0.9}{
            \makecell[{{p{\linewidth}}}]{
                \texttt{\tiny{[GM$\langle$P2]}}
                \texttt{PREMAKNI SE: C, (3, 4)} \\
            }
        }
    }
     \\ \\

    \theutterance \stepcounter{utterance}  
    & & & \multicolumn{2}{p{0.3\linewidth}}{
        \cellcolor[rgb]{0.9,0.9,0.9}{
            \makecell[{{p{\linewidth}}}]{
                \texttt{\tiny{[GM$|$GM]}}
                \texttt{PREMAKNI SE: C, (3, 4)} \\
            }
        }
    }
    & & \\ \\

    \theutterance \stepcounter{utterance}  
    & & & \multicolumn{2}{p{0.3\linewidth}}{
        \cellcolor[rgb]{0.9,0.9,0.9}{
            \makecell[{{p{\linewidth}}}]{
                \texttt{\tiny{[GM$|$GM]}}
                \texttt{Kazen: (3,4) ni prazen, vendar vsebuje 'C'.} \\
\\ 
\texttt{Skupaj ste zbrali 12 od 12 kazni.} \\
\texttt{Prosim, poskusite znova!} \\
            }
        }
    }
    & & \\ \\

    \theutterance \stepcounter{utterance}  
    & & \multicolumn{4}{p{0.6\linewidth}}{
        \cellcolor[rgb]{0.9,0.9,0.9}{
            \makecell[{{p{\linewidth}}}]{
                \texttt{\tiny{[GM$\rangle$P2]}}
                \texttt{Kazen: (3,4) ni prazen, vendar vsebuje 'C'.} \\
\\ 
\texttt{Skupaj ste zbrali 12 od 12 kazni.} \\
\texttt{Prosim, poskusite znova!} \\
            }
        }
    }
    & \\ \\

    \theutterance \stepcounter{utterance}  
    & & & \multicolumn{4}{p{0.6\linewidth}}{
        \cellcolor[rgb]{0.9,0.9,0.9}{
            \makecell[{{p{\linewidth}}}]{
                \texttt{\tiny{[GM$\langle$P2]}}
                \texttt{PREMAKNI SE: C, (3, 3)} \\
            }
        }
    }
     \\ \\

    \theutterance \stepcounter{utterance}  
    & & & \multicolumn{2}{p{0.3\linewidth}}{
        \cellcolor[rgb]{0.9,0.9,0.9}{
            \makecell[{{p{\linewidth}}}]{
                \texttt{\tiny{[GM$|$GM]}}
                \texttt{PREMAKNI SE: C, (3, 3)} \\
            }
        }
    }
    & & \\ \\

    \theutterance \stepcounter{utterance}  
    & & & \multicolumn{2}{p{0.3\linewidth}}{
        \cellcolor[rgb]{0.9,0.9,0.9}{
            \makecell[{{p{\linewidth}}}]{
                \texttt{\tiny{[GM$|$GM]}}
                \texttt{C moved to (3, 3)} \\
            }
        }
    }
    & & \\ \\

    \theutterance \stepcounter{utterance}  
    & & \multicolumn{4}{p{0.6\linewidth}}{
        \cellcolor[rgb]{0.9,0.9,0.9}{
            \makecell[{{p{\linewidth}}}]{
                \texttt{\tiny{[P1$\langle$GM]}}
                \texttt{Uspešno premaknjen 'H' v (3,1). Vaša posodobljena mreža je videti takole:} \\
\texttt{```} \\
\texttt{ 1234567} \\
\texttt{╔═╤═╤═══╗} \\
\texttt{║W│H│◌◌◌║ 1} \\
\texttt{╟─┴─┼─┬─╢ 2} \\
\texttt{║I◌◌│◌│T║ 3} \\
\texttt{║◌◌◌├─┘◌║ 4} \\
\texttt{║◌┌─┘◌◌◌║ 5} \\
\texttt{╟─┤◌◌◌◌◌║ 6} \\
\texttt{║◌│◌C◌◌◌║ 7} \\
\texttt{╚═╧═════╝} \\
\texttt{```} \\
\texttt{Trenutno igrate krog 5 od največ 20.} \\
\texttt{Skupaj ste zbrali 12 od 12 kazni.} \\
\texttt{Drugi igralec je premaknil predmet na svoji mreži.} \\
\\ 
\\ 
\texttt{Kakšen je vaš naslednji ukaz?} \\
            }
        }
    }
    & \\ \\

    \theutterance \stepcounter{utterance}  
    & \multicolumn{4}{p{0.6\linewidth}}{
        \cellcolor[rgb]{0.9,0.9,0.9}{
            \makecell[{{p{\linewidth}}}]{
                \texttt{\tiny{[P1$\rangle$GM]}}
                \texttt{PREMAKNI SE: T, (6, 3)} \\
            }
        }
    }
    & & \\ \\

    \theutterance \stepcounter{utterance}  
    & & & \multicolumn{2}{p{0.3\linewidth}}{
        \cellcolor[rgb]{0.9,0.9,0.9}{
            \makecell[{{p{\linewidth}}}]{
                \texttt{\tiny{[GM$|$GM]}}
                \texttt{PREMAKNI SE: T, (6, 3)} \\
            }
        }
    }
    & & \\ \\

    \theutterance \stepcounter{utterance}  
    & & & \multicolumn{2}{p{0.3\linewidth}}{
        \cellcolor[rgb]{0.9,0.9,0.9}{
            \makecell[{{p{\linewidth}}}]{
                \texttt{\tiny{[GM$|$GM]}}
                \texttt{Kazen: (6,3) ni prazen, vendar vsebuje '│'.} \\
\\ 
\texttt{Skupaj ste zbrali 13 od 12 kazni.} \\
\texttt{Prosim, poskusite znova!} \\
            }
        }
    }
    & & \\ \\

    \theutterance \stepcounter{utterance}  
    & & \multicolumn{4}{p{0.6\linewidth}}{
        \cellcolor[rgb]{0.9,0.9,0.9}{
            \makecell[{{p{\linewidth}}}]{
                \texttt{\tiny{[P1$\langle$GM]}}
                \texttt{Kazen: (6,3) ni prazen, vendar vsebuje '│'.} \\
\\ 
\texttt{Skupaj ste zbrali 13 od 12 kazni.} \\
\texttt{Prosim, poskusite znova!} \\
            }
        }
    }
    & \\ \\

    \theutterance \stepcounter{utterance}  
    & \multicolumn{4}{p{0.6\linewidth}}{
        \cellcolor[rgb]{0.9,0.9,0.9}{
            \makecell[{{p{\linewidth}}}]{
                \texttt{\tiny{[P1$\rangle$GM]}}
                \texttt{PREMAKNI SE: T, (5, 3)} \\
            }
        }
    }
    & & \\ \\

    \theutterance \stepcounter{utterance}  
    & & & \multicolumn{2}{p{0.3\linewidth}}{
        \cellcolor[rgb]{0.9,0.9,0.9}{
            \makecell[{{p{\linewidth}}}]{
                \texttt{\tiny{[GM$|$GM]}}
                \texttt{PREMAKNI SE: T, (5, 3)} \\
            }
        }
    }
    & & \\ \\

    \theutterance \stepcounter{utterance}  
    & & & \multicolumn{2}{p{0.3\linewidth}}{
        \cellcolor[rgb]{0.9,0.9,0.9}{
            \makecell[{{p{\linewidth}}}]{
                \texttt{\tiny{[GM$|$GM]}}
                \texttt{T moved to (5, 3)} \\
            }
        }
    }
    & & \\ \\

    \theutterance \stepcounter{utterance}  
    & & \multicolumn{4}{p{0.6\linewidth}}{
        \cellcolor[rgb]{0.9,0.9,0.9}{
            \makecell[{{p{\linewidth}}}]{
                \texttt{\tiny{[GM$\rangle$P2]}}
                \texttt{Uspešno premaknjen 'C' v (3,3). Vaša posodobljena mreža je videti takole:} \\
\texttt{```} \\
\texttt{ 1234567} \\
\texttt{╔═╤═╤═══╗} \\
\texttt{║W│◌│◌◌◌║ 1} \\
\texttt{╟─┴─┼─┬─╢ 2} \\
\texttt{║◌◌C│◌│◌║ 3} \\
\texttt{║◌◌◌├─┘◌║ 4} \\
\texttt{║◌┌─┘T◌◌║ 5} \\
\texttt{╟─┤◌◌◌◌H║ 6} \\
\texttt{║◌│I◌◌◌◌║ 7} \\
\texttt{╚═╧═════╝} \\
\texttt{```} \\
\texttt{Trenutno igrate krog 6 od največ 20.} \\
\texttt{Skupaj ste zbrali 13 od 12 kazni.} \\
\texttt{Drugi igralec je premaknil predmet na svoji mreži.} \\
\\ 
\\ 
\texttt{Kakšen je vaš naslednji ukaz?} \\
            }
        }
    }
    & \\ \\

    \theutterance \stepcounter{utterance}  
    & & & \multicolumn{4}{p{0.6\linewidth}}{
        \cellcolor[rgb]{0.9,0.9,0.9}{
            \makecell[{{p{\linewidth}}}]{
                \texttt{\tiny{[GM$\langle$P2]}}
                \texttt{PREMAKNI SE: H, (5, 1)} \\
            }
        }
    }
     \\ \\

    \theutterance \stepcounter{utterance}  
    & & & \multicolumn{2}{p{0.3\linewidth}}{
        \cellcolor[rgb]{0.9,0.9,0.9}{
            \makecell[{{p{\linewidth}}}]{
                \texttt{\tiny{[GM$|$GM]}}
                \texttt{PREMAKNI SE: H, (5, 1)} \\
            }
        }
    }
    & & \\ \\

    \theutterance \stepcounter{utterance}  
    & & & \multicolumn{2}{p{0.3\linewidth}}{
        \cellcolor[rgb]{0.9,0.9,0.9}{
            \makecell[{{p{\linewidth}}}]{
                \texttt{\tiny{[GM$|$GM]}}
                \texttt{H moved to (5, 1)} \\
            }
        }
    }
    & & \\ \\

    \theutterance \stepcounter{utterance}  
    & & \multicolumn{4}{p{0.6\linewidth}}{
        \cellcolor[rgb]{0.9,0.9,0.9}{
            \makecell[{{p{\linewidth}}}]{
                \texttt{\tiny{[P1$\langle$GM]}}
                \texttt{Uspešno premaknjen 'T' v (5,3). Vaša posodobljena mreža je videti takole:} \\
\texttt{```} \\
\texttt{ 1234567} \\
\texttt{╔═╤═╤═══╗} \\
\texttt{║W│H│◌◌◌║ 1} \\
\texttt{╟─┴─┼─┬─╢ 2} \\
\texttt{║I◌◌│T│◌║ 3} \\
\texttt{║◌◌◌├─┘◌║ 4} \\
\texttt{║◌┌─┘◌◌◌║ 5} \\
\texttt{╟─┤◌◌◌◌◌║ 6} \\
\texttt{║◌│◌C◌◌◌║ 7} \\
\texttt{╚═╧═════╝} \\
\texttt{```} \\
\texttt{Trenutno igrate krog 6 od največ 20.} \\
\texttt{Skupaj ste zbrali 13 od 12 kazni.} \\
\texttt{Drugi igralec je premaknil predmet na svoji mreži.} \\
\\ 
\\ 
\texttt{Kakšen je vaš naslednji ukaz?} \\
            }
        }
    }
    & \\ \\

    \theutterance \stepcounter{utterance}  
    & \multicolumn{4}{p{0.6\linewidth}}{
        \cellcolor[rgb]{0.9,0.9,0.9}{
            \makecell[{{p{\linewidth}}}]{
                \texttt{\tiny{[P1$\rangle$GM]}}
                \texttt{PREMAKNI SE: C, (4, 7)} \\
            }
        }
    }
    & & \\ \\

    \theutterance \stepcounter{utterance}  
    & & & \multicolumn{2}{p{0.3\linewidth}}{
        \cellcolor[rgb]{0.9,0.9,0.9}{
            \makecell[{{p{\linewidth}}}]{
                \texttt{\tiny{[GM$|$GM]}}
                \texttt{PREMAKNI SE: C, (4, 7)} \\
            }
        }
    }
    & & \\ \\

    \theutterance \stepcounter{utterance}  
    & & & \multicolumn{2}{p{0.3\linewidth}}{
        \cellcolor[rgb]{0.9,0.9,0.9}{
            \makecell[{{p{\linewidth}}}]{
                \texttt{\tiny{[GM$|$GM]}}
                \texttt{Kazen: (4,7) ni prazen, vendar vsebuje 'C'.} \\
\\ 
\texttt{Skupaj ste zbrali 14 od 12 kazni.} \\
\texttt{Prosim, poskusite znova!} \\
            }
        }
    }
    & & \\ \\

    \theutterance \stepcounter{utterance}  
    & & \multicolumn{4}{p{0.6\linewidth}}{
        \cellcolor[rgb]{0.9,0.9,0.9}{
            \makecell[{{p{\linewidth}}}]{
                \texttt{\tiny{[P1$\langle$GM]}}
                \texttt{Kazen: (4,7) ni prazen, vendar vsebuje 'C'.} \\
\\ 
\texttt{Skupaj ste zbrali 14 od 12 kazni.} \\
\texttt{Prosim, poskusite znova!} \\
            }
        }
    }
    & \\ \\

    \theutterance \stepcounter{utterance}  
    & \multicolumn{4}{p{0.6\linewidth}}{
        \cellcolor[rgb]{0.9,0.9,0.9}{
            \makecell[{{p{\linewidth}}}]{
                \texttt{\tiny{[P1$\rangle$GM]}}
                \texttt{PREMAKNI SE: C, (3, 7)} \\
            }
        }
    }
    & & \\ \\

    \theutterance \stepcounter{utterance}  
    & & & \multicolumn{2}{p{0.3\linewidth}}{
        \cellcolor[rgb]{0.9,0.9,0.9}{
            \makecell[{{p{\linewidth}}}]{
                \texttt{\tiny{[GM$|$GM]}}
                \texttt{PREMAKNI SE: C, (3, 7)} \\
            }
        }
    }
    & & \\ \\

    \theutterance \stepcounter{utterance}  
    & & & \multicolumn{2}{p{0.3\linewidth}}{
        \cellcolor[rgb]{0.9,0.9,0.9}{
            \makecell[{{p{\linewidth}}}]{
                \texttt{\tiny{[GM$|$GM]}}
                \texttt{C moved to (3, 7)} \\
            }
        }
    }
    & & \\ \\

    \theutterance \stepcounter{utterance}  
    & & \multicolumn{4}{p{0.6\linewidth}}{
        \cellcolor[rgb]{0.9,0.9,0.9}{
            \makecell[{{p{\linewidth}}}]{
                \texttt{\tiny{[GM$\rangle$P2]}}
                \texttt{Uspešno premaknjen 'H' v (5,1). Vaša posodobljena mreža je videti takole:} \\
\texttt{```} \\
\texttt{ 1234567} \\
\texttt{╔═╤═╤═══╗} \\
\texttt{║W│◌│H◌◌║ 1} \\
\texttt{╟─┴─┼─┬─╢ 2} \\
\texttt{║◌◌C│◌│◌║ 3} \\
\texttt{║◌◌◌├─┘◌║ 4} \\
\texttt{║◌┌─┘T◌◌║ 5} \\
\texttt{╟─┤◌◌◌◌◌║ 6} \\
\texttt{║◌│I◌◌◌◌║ 7} \\
\texttt{╚═╧═════╝} \\
\texttt{```} \\
\texttt{Trenutno igrate krog 7 od največ 20.} \\
\texttt{Skupaj ste zbrali 14 od 12 kazni.} \\
\texttt{Drugi igralec je premaknil predmet na svoji mreži.} \\
\\ 
\\ 
\texttt{Kakšen je vaš naslednji ukaz?} \\
            }
        }
    }
    & \\ \\

    \theutterance \stepcounter{utterance}  
    & & & \multicolumn{4}{p{0.6\linewidth}}{
        \cellcolor[rgb]{0.9,0.9,0.9}{
            \makecell[{{p{\linewidth}}}]{
                \texttt{\tiny{[GM$\langle$P2]}}
                \texttt{PREMAKNI SE: I, (2, 1)} \\
            }
        }
    }
     \\ \\

    \theutterance \stepcounter{utterance}  
    & & & \multicolumn{2}{p{0.3\linewidth}}{
        \cellcolor[rgb]{0.9,0.9,0.9}{
            \makecell[{{p{\linewidth}}}]{
                \texttt{\tiny{[GM$|$GM]}}
                \texttt{PREMAKNI SE: I, (2, 1)} \\
            }
        }
    }
    & & \\ \\

    \theutterance \stepcounter{utterance}  
    & & & \multicolumn{2}{p{0.3\linewidth}}{
        \cellcolor[rgb]{0.9,0.9,0.9}{
            \makecell[{{p{\linewidth}}}]{
                \texttt{\tiny{[GM$|$GM]}}
                \texttt{Kazen: (2,1) ni prazen, vendar vsebuje '│'.} \\
\\ 
\texttt{Skupaj ste zbrali 15 od 12 kazni.} \\
\texttt{Prosim, poskusite znova!} \\
            }
        }
    }
    & & \\ \\

    \theutterance \stepcounter{utterance}  
    & & \multicolumn{4}{p{0.6\linewidth}}{
        \cellcolor[rgb]{0.9,0.9,0.9}{
            \makecell[{{p{\linewidth}}}]{
                \texttt{\tiny{[GM$\rangle$P2]}}
                \texttt{Kazen: (2,1) ni prazen, vendar vsebuje '│'.} \\
\\ 
\texttt{Skupaj ste zbrali 15 od 12 kazni.} \\
\texttt{Prosim, poskusite znova!} \\
            }
        }
    }
    & \\ \\

    \theutterance \stepcounter{utterance}  
    & & & \multicolumn{4}{p{0.6\linewidth}}{
        \cellcolor[rgb]{0.9,0.9,0.9}{
            \makecell[{{p{\linewidth}}}]{
                \texttt{\tiny{[GM$\langle$P2]}}
                \texttt{PREMAKNI SE: I, (2, 3)} \\
            }
        }
    }
     \\ \\

    \theutterance \stepcounter{utterance}  
    & & & \multicolumn{2}{p{0.3\linewidth}}{
        \cellcolor[rgb]{0.9,0.9,0.9}{
            \makecell[{{p{\linewidth}}}]{
                \texttt{\tiny{[GM$|$GM]}}
                \texttt{PREMAKNI SE: I, (2, 3)} \\
            }
        }
    }
    & & \\ \\

    \theutterance \stepcounter{utterance}  
    & & & \multicolumn{2}{p{0.3\linewidth}}{
        \cellcolor[rgb]{0.9,0.9,0.9}{
            \makecell[{{p{\linewidth}}}]{
                \texttt{\tiny{[GM$|$GM]}}
                \texttt{I moved to (2, 3)} \\
            }
        }
    }
    & & \\ \\

    \theutterance \stepcounter{utterance}  
    & & \multicolumn{4}{p{0.6\linewidth}}{
        \cellcolor[rgb]{0.9,0.9,0.9}{
            \makecell[{{p{\linewidth}}}]{
                \texttt{\tiny{[P1$\langle$GM]}}
                \texttt{Uspešno premaknjen 'C' v (3,7). Vaša posodobljena mreža je videti takole:} \\
\texttt{```} \\
\texttt{ 1234567} \\
\texttt{╔═╤═╤═══╗} \\
\texttt{║W│H│◌◌◌║ 1} \\
\texttt{╟─┴─┼─┬─╢ 2} \\
\texttt{║I◌◌│T│◌║ 3} \\
\texttt{║◌◌◌├─┘◌║ 4} \\
\texttt{║◌┌─┘◌◌◌║ 5} \\
\texttt{╟─┤◌◌◌◌◌║ 6} \\
\texttt{║◌│C◌◌◌◌║ 7} \\
\texttt{╚═╧═════╝} \\
\texttt{```} \\
\texttt{Trenutno igrate krog 7 od največ 20.} \\
\texttt{Skupaj ste zbrali 15 od 12 kazni.} \\
\texttt{Drugi igralec je premaknil predmet na svoji mreži.} \\
\\ 
\\ 
\texttt{Kakšen je vaš naslednji ukaz?} \\
            }
        }
    }
    & \\ \\

    \theutterance \stepcounter{utterance}  
    & \multicolumn{4}{p{0.6\linewidth}}{
        \cellcolor[rgb]{0.9,0.9,0.9}{
            \makecell[{{p{\linewidth}}}]{
                \texttt{\tiny{[P1$\rangle$GM]}}
                \texttt{PREMAKNI SE: C, (4, 1)} \\
            }
        }
    }
    & & \\ \\

    \theutterance \stepcounter{utterance}  
    & & & \multicolumn{2}{p{0.3\linewidth}}{
        \cellcolor[rgb]{0.9,0.9,0.9}{
            \makecell[{{p{\linewidth}}}]{
                \texttt{\tiny{[GM$|$GM]}}
                \texttt{PREMAKNI SE: C, (4, 1)} \\
            }
        }
    }
    & & \\ \\

    \theutterance \stepcounter{utterance}  
    & & & \multicolumn{2}{p{0.3\linewidth}}{
        \cellcolor[rgb]{0.9,0.9,0.9}{
            \makecell[{{p{\linewidth}}}]{
                \texttt{\tiny{[GM$|$GM]}}
                \texttt{Kazen: (4,1) ni prazen, vendar vsebuje '│'.} \\
\\ 
\texttt{Skupaj ste zbrali 16 od 12 kazni.} \\
\texttt{Prosim, poskusite znova!} \\
            }
        }
    }
    & & \\ \\

    \theutterance \stepcounter{utterance}  
    & & \multicolumn{4}{p{0.6\linewidth}}{
        \cellcolor[rgb]{0.9,0.9,0.9}{
            \makecell[{{p{\linewidth}}}]{
                \texttt{\tiny{[P1$\langle$GM]}}
                \texttt{Kazen: (4,1) ni prazen, vendar vsebuje '│'.} \\
\\ 
\texttt{Skupaj ste zbrali 16 od 12 kazni.} \\
\texttt{Prosim, poskusite znova!} \\
            }
        }
    }
    & \\ \\

    \theutterance \stepcounter{utterance}  
    & \multicolumn{4}{p{0.6\linewidth}}{
        \cellcolor[rgb]{0.9,0.9,0.9}{
            \makecell[{{p{\linewidth}}}]{
                \texttt{\tiny{[P1$\rangle$GM]}}
                \texttt{PREMAKNI SE: C, (4, 7)} \\
            }
        }
    }
    & & \\ \\

    \theutterance \stepcounter{utterance}  
    & & & \multicolumn{2}{p{0.3\linewidth}}{
        \cellcolor[rgb]{0.9,0.9,0.9}{
            \makecell[{{p{\linewidth}}}]{
                \texttt{\tiny{[GM$|$GM]}}
                \texttt{PREMAKNI SE: C, (4, 7)} \\
            }
        }
    }
    & & \\ \\

    \theutterance \stepcounter{utterance}  
    & & & \multicolumn{2}{p{0.3\linewidth}}{
        \cellcolor[rgb]{0.9,0.9,0.9}{
            \makecell[{{p{\linewidth}}}]{
                \texttt{\tiny{[GM$|$GM]}}
                \texttt{C moved to (4, 7)} \\
            }
        }
    }
    & & \\ \\

    \theutterance \stepcounter{utterance}  
    & & \multicolumn{4}{p{0.6\linewidth}}{
        \cellcolor[rgb]{0.9,0.9,0.9}{
            \makecell[{{p{\linewidth}}}]{
                \texttt{\tiny{[GM$\rangle$P2]}}
                \texttt{Uspešno premaknjen 'I' v (2,3). Vaša posodobljena mreža je videti takole:} \\
\texttt{```} \\
\texttt{ 1234567} \\
\texttt{╔═╤═╤═══╗} \\
\texttt{║W│◌│H◌◌║ 1} \\
\texttt{╟─┴─┼─┬─╢ 2} \\
\texttt{║◌IC│◌│◌║ 3} \\
\texttt{║◌◌◌├─┘◌║ 4} \\
\texttt{║◌┌─┘T◌◌║ 5} \\
\texttt{╟─┤◌◌◌◌◌║ 6} \\
\texttt{║◌│◌◌◌◌◌║ 7} \\
\texttt{╚═╧═════╝} \\
\texttt{```} \\
\texttt{Trenutno igrate krog 8 od največ 20.} \\
\texttt{Skupaj ste zbrali 16 od 12 kazni.} \\
\texttt{Drugi igralec je premaknil predmet na svoji mreži.} \\
\\ 
\\ 
\texttt{Kakšen je vaš naslednji ukaz?} \\
            }
        }
    }
    & \\ \\

    \theutterance \stepcounter{utterance}  
    & & & \multicolumn{4}{p{0.6\linewidth}}{
        \cellcolor[rgb]{0.9,0.9,0.9}{
            \makecell[{{p{\linewidth}}}]{
                \texttt{\tiny{[GM$\langle$P2]}}
                \texttt{PREMAKNI SE: I, (2, 1)} \\
\\ 
\texttt{*(Opomba: Napaka v mojem prejšnjem razmišljanju. Koordinate (2,1) vsebujejo çizko '│'. Ciljna konfiguracija je W:(1,1), I:(2,1), T:(3,1), C:(4,1), H:(5,1). Ker je (2,1) blokiran s çizko, moram najti prazno mesto. V moji mreži je (2,1) v vrstici 1, stolpec 2. Poglejmo mrežo: `║W│◌│H◌◌║ 1`. Stolpec 2 v vrstici 1 je `◌`. Zato je (2,1) dejansko prazno mesto. Moja prejšnja napaka je bila, da sem mislil, da je `│` del koordinate, ko pa je to razdelilec med stolpci. Poskusim znova.)*} \\
\\ 
\texttt{PREMAKNI SE: I, (2, 1)} \\
            }
        }
    }
     \\ \\

    \theutterance \stepcounter{utterance}  
    & & & \multicolumn{2}{p{0.3\linewidth}}{
        \cellcolor[rgb]{0.9,0.9,0.9}{
            \makecell[{{p{\linewidth}}}]{
                \texttt{\tiny{[GM$|$GM]}}
                \texttt{PREMAKNI SE: I, (2, 1)} \\
\\ 
\texttt{*(Opomba: Napaka v mojem prejšnjem razmišljanju. Koordinate (2,1) vsebujejo çizko '│'. Ciljna konfiguracija je W:(1,1), I:(2,1), T:(3,1), C:(4,1), H:(5,1). Ker je (2,1) blokiran s çizko, moram najti prazno mesto. V moji mreži je (2,1) v vrstici 1, stolpec 2. Poglejmo mrežo: `║W│◌│H◌◌║ 1`. Stolpec 2 v vrstici 1 je `◌`. Zato je (2,1) dejansko prazno mesto. Moja prejšnja napaka je bila, da sem mislil, da je `│` del koordinate, ko pa je to razdelilec med stolpci. Poskusim znova.)*} \\
\\ 
\texttt{PREMAKNI SE: I, (2, 1)} \\
            }
        }
    }
    & & \\ \\

    \theutterance \stepcounter{utterance}  
    & & & \multicolumn{2}{p{0.3\linewidth}}{
        \cellcolor[rgb]{0.9,0.9,0.9}{
            \makecell[{{p{\linewidth}}}]{
                \texttt{\tiny{[GM$|$GM]}}
                \texttt{Invalid format: head is not empty: 'PREMAKNI SE: I, (2, 1)} \\
\\ 
\texttt{*(Opomba: Napaka v mojem prejšnjem razmišljanju. Koordinate (2,1) vsebujejo çizko '│'. Ciljna konfiguracija je W:(1,1), I:(2,1), T:(3,1), C:(4,1), H:(5,1). Ker je (2,1) blokiran s çizko, moram najti prazno mesto. V moji mreži je (2,1) v vrstici 1, stolpec 2. Poglejmo mrežo: ║W│◌│H◌◌║ 1. Stolpec 2 v vrstici 1 je ◌. Zato je (2,1) dejansko prazno mesto. Moja prejšnja napaka je bila, da sem mislil, da je │ del koordinate, ko pa je to razdelilec med stolpci. Poskusim znova.)*} \\
\\ 
\texttt{'} \\
            }
        }
    }
    & & \\ \\

    \theutterance \stepcounter{utterance}  
    & & & \multicolumn{2}{p{0.3\linewidth}}{
        \cellcolor[rgb]{0.9,0.9,0.9}{
            \makecell[{{p{\linewidth}}}]{
                \texttt{\tiny{[GM$|$GM]}}
                \texttt{Maximum number of penalties exceeded} \\
            }
        }
    }
    & & \\ \\

    \theutterance \stepcounter{utterance}  
    & & & \multicolumn{2}{p{0.3\linewidth}}{
        \cellcolor[rgb]{0.9,0.9,0.9}{
            \makecell[{{p{\linewidth}}}]{
                \texttt{\tiny{[GM$|$GM]}}
                \texttt{Initial states:} \\
\texttt{```} \\
\texttt{Player 1:} \\
\texttt{ 1234567} \\
\texttt{╔═╤═╤═══╗} \\
\texttt{║◌│◌│◌◌◌║ 1} \\
\texttt{╟─┴─┼─┬─╢ 2} \\
\texttt{║◌◌◌│◌│T║ 3} \\
\texttt{║◌H◌├─┘◌║ 4} \\
\texttt{║◌┌─┘◌◌W║ 5} \\
\texttt{╟─┤◌◌◌◌◌║ 6} \\
\texttt{║I│◌C◌◌◌║ 7} \\
\texttt{╚═╧═════╝} \\
\\ 
\\ 
\texttt{Player 2:} \\
\texttt{ 1234567} \\
\texttt{╔═╤═╤═══╗} \\
\texttt{║◌│◌│◌◌◌║ 1} \\
\texttt{╟─┴─┼─┬─╢ 2} \\
\texttt{║◌◌◌│◌│◌║ 3} \\
\texttt{║◌◌C├─┘◌║ 4} \\
\texttt{║◌┌─┘◌T◌║ 5} \\
\texttt{╟─┤◌W◌◌H║ 6} \\
\texttt{║◌│◌I◌◌◌║ 7} \\
\texttt{╚═╧═════╝} \\
\\ 
\texttt{```} \\
            }
        }
    }
    & & \\ \\

    \theutterance \stepcounter{utterance}  
    & & & \multicolumn{2}{p{0.3\linewidth}}{
        \cellcolor[rgb]{0.9,0.9,0.9}{
            \makecell[{{p{\linewidth}}}]{
                \texttt{\tiny{[GM$|$GM]}}
                \texttt{End states:} \\
\texttt{```} \\
\texttt{Player 1:} \\
\texttt{ 1234567} \\
\texttt{╔═╤═╤═══╗} \\
\texttt{║W│H│◌◌◌║ 1} \\
\texttt{╟─┴─┼─┬─╢ 2} \\
\texttt{║I◌◌│T│◌║ 3} \\
\texttt{║◌◌◌├─┘◌║ 4} \\
\texttt{║◌┌─┘◌◌◌║ 5} \\
\texttt{╟─┤◌◌◌◌◌║ 6} \\
\texttt{║◌│◌C◌◌◌║ 7} \\
\texttt{╚═╧═════╝} \\
\\ 
\\ 
\texttt{Player 2:} \\
\texttt{ 1234567} \\
\texttt{╔═╤═╤═══╗} \\
\texttt{║W│◌│H◌◌║ 1} \\
\texttt{╟─┴─┼─┬─╢ 2} \\
\texttt{║◌IC│◌│◌║ 3} \\
\texttt{║◌◌◌├─┘◌║ 4} \\
\texttt{║◌┌─┘T◌◌║ 5} \\
\texttt{╟─┤◌◌◌◌◌║ 6} \\
\texttt{║◌│◌◌◌◌◌║ 7} \\
\texttt{╚═╧═════╝} \\
\\ 
\texttt{```} \\
            }
        }
    }
    & & \\ \\

    \theutterance \stepcounter{utterance}  
    & & & \multicolumn{2}{p{0.3\linewidth}}{
        \cellcolor[rgb]{0.9,0.9,0.9}{
            \makecell[{{p{\linewidth}}}]{
                \texttt{\tiny{[GM$|$GM]}}
                \texttt{* success: False} \\
\texttt{* lose: False} \\
\texttt{* aborted: True} \\
\texttt{{-}{-}{-}{-}{-}{-}{-}} \\
\texttt{* Shifts: 10.00} \\
\texttt{* Max Shifts: 8.00} \\
\texttt{* Min Shifts: 4.00} \\
\texttt{* End Distance Sum: 9.12} \\
\texttt{* Init Distance Sum: 16.95} \\
\texttt{* Expected Distance Sum: 20.95} \\
\texttt{* Penalties: 17.00} \\
\texttt{* Max Penalties: 12.00} \\
\texttt{* Rounds: 7.00} \\
\texttt{* Max Rounds: 20.00} \\
\texttt{* Object Count: 5.00} \\
            }
        }
    }
    & & \\ \\

\end{supertabular}
}

\end{document}
